% Imporditud: tools/import_eksperimendid.py
% Allikas: Eesti füüsikaolümpiaadi eksperimendi ülesannete
%   kogu (Rael Kalda, 2023), ülesanne Ü17, lahendus L17.
\setAuthor{EFO žürii}
\setRound{piirkonnavoor}
\setYear{2009}
\setNumber{P E1}
\setDifficulty{1}
\setTopic{vedelike mehaanika}
\setVahendid{mõõtemensuur, anum veega, haavlid}
\setPunktid{10}
\setAllikas{Eksperimendi kogu Ü17, lahendus lk 39}
\setLahendusseis{toores}

\prob{Haavlid}
Millise suhtelise osa haavlitega (sh tühikud) täidetud ruumalast moodustavad tühikud haavlite vahel?

\hint

\solu
Kallame mensuuri vett ja fikseerime skaalal vedeliku mahu näidu $V_0$. Lisame seejärel haavleid, niiet veetase mensuuris jääks skaala ulatusse ja ülespoole haavlite taset (kerge loksutamisega sätime viimase horisontaalseks). Fikseerime koguruumala $V_1$, mis jääb allapoole veetaset ja haavlite + veega täidetud osa ruumala $V_2$. Veehulga jäävusest enne ja pärast haavlite lisamist

$V_0$ = $V_2$ $\cdot$ S + ($V_1$ $-$ $V_2$),

kus S on otsitav tühikute suhteline ruumala. Sellest seosest saame

S = 1 $-$ ($V_1$ $-$ $V_0$)/$V_2$.

Erijuhul kui haavlite tase ühtib veetasemega ($V_1$ = $V_2$), siis S = $V_0$/$V_2$.

Märkus: Sama läbimõõduga kerakeste juhusliku pakkimise korral on näidatud, et S = 0,38. Tihedaima (regulaarse) pakendi korral S = 0,26.

Hindamine: Kirjeldatud õige mõõtmisprotseduur ja arvutuseeskiri --- 4 p. Teostatud mõõtmised ja esitatud tulemused --- 2 p. Arvutatud S väärtus --- 2 p. Tühikutega täidetud osa ruumala vahemikus 0,25 $\leq$ S $\leq$ 0,50 --- 2 p.
\probend
